\documentclass{article}
\usepackage{fuzz}

\begin{document}

\title{vox-2594: Concurrent-Writer Log Rotation}
\author{rmh}
\maketitle

% The processes that append to vox.log: the daemon and every transient client
% (hook, MCP server, CLI, detached playback). Each takes a shared flock to
% append and upgrades to an exclusive flock to rotate.
\begin{zed}
  [WRITER]
\end{zed}

% The lock state of one vox.log sink guarded by a stable rotate.lock file.
%   shared     -- writers holding LOCK_SH, each mid open/write/close append.
%   exclusive  -- the writer holding LOCK_EX to run the rename chain.
%   size       -- current bytes in the active vox.log.
%   backups    -- how many rotated backup slots exist (.1 .. .backupCount).
%   maxBytes   -- rotation threshold.
%   backupCount-- retained backup slots.
%
% The predicate carries the safety invariants the flock protocol must hold:
%   (I1) a writer cannot hold both a shared and the exclusive lock at once;
%   (I2) at most one rotator runs the rename chain;
%   (I3) a rotation excludes every appender -- no append fd is open while a
%        rename runs, so no line is written to a renamed inode;
%   (I4) the backup chain never exceeds its retained length.
\begin{schema}{Rotator}
  shared: \power WRITER \\
  exclusive: \power WRITER \\
  size: \nat \\
  backups: \nat \\
  maxBytes: \nat \\
  backupCount: \nat
\where
  shared \cap exclusive = \emptyset \\
  \# exclusive \leq 1 \\
  exclusive \neq \emptyset \implies shared = \emptyset \\
  backups \leq backupCount
\end{schema}

% Initially no lock is held, the log is empty, and no backups exist. The
% thresholds are positive so a rotation is reachable and a backup can be kept.
\begin{schema}{InitRotator}
  Rotator
\where
  shared = \emptyset \\
  exclusive = \emptyset \\
  size = 0 \\
  backups = 0 \\
  maxBytes > 0 \\
  backupCount > 0
\end{schema}

% A writer takes LOCK_SH to begin an append. Refused while any rotator holds
% LOCK_EX (I3): flock blocks a shared acquire behind an exclusive holder.
\begin{schema}{AcquireShared}
  \Delta Rotator \\
  w?: WRITER
\where
  exclusive = \emptyset \\
  w? \notin shared \\
  shared' = shared \cup \{ w? \} \\
  exclusive' = exclusive \\
  size' = size \\
  backups' = backups \\
  maxBytes' = maxBytes \\
  backupCount' = backupCount
\end{schema}

% The atomic O_APPEND write of one line, performed while holding LOCK_SH. Only
% fires when the line fits under the threshold; an oversize line rotates first.
\begin{schema}{Append}
  \Delta Rotator \\
  w?: WRITER \\
  len?: \nat
\where
  w? \in shared \\
  size + len? \leq maxBytes \\
  size' = size + len? \\
  shared' = shared \\
  exclusive' = exclusive \\
  backups' = backups \\
  maxBytes' = maxBytes \\
  backupCount' = backupCount
\end{schema}

% The writer releases LOCK_SH after close, allowing a pending rotator to proceed.
\begin{schema}{ReleaseShared}
  \Delta Rotator \\
  w?: WRITER
\where
  w? \in shared \\
  shared' = shared \setminus \{ w? \} \\
  exclusive' = exclusive \\
  size' = size \\
  backups' = backups \\
  maxBytes' = maxBytes \\
  backupCount' = backupCount
\end{schema}

% A writer that observed an oversize append upgrades to LOCK_EX. It has already
% released its own shared lock, and flock grants the exclusive lock only once
% every shared holder has drained (shared = empty) and no other rotator holds it
% (exclusive = empty). This establishes I3 for the rename that follows.
\begin{schema}{BeginRotate}
  \Delta Rotator \\
  w?: WRITER
\where
  w? \notin shared \\
  shared = \emptyset \\
  exclusive = \emptyset \\
  exclusive' = \{ w? \} \\
  shared' = shared \\
  size' = size \\
  backups' = backups \\
  maxBytes' = maxBytes \\
  backupCount' = backupCount
\end{schema}

% Under LOCK_EX, re-check found the file still oversize: run the rename chain.
% The active file resets to empty and one more backup slot exists, capped at
% backupCount (the oldest backup is overwritten, not accumulated).
\begin{schema}{Rotate}
  \Delta Rotator \\
  w?: WRITER
\where
  w? \in exclusive \\
  size > maxBytes \\
  size' = 0 \\
  backups < backupCount \implies backups' = backups + 1 \\
  backups = backupCount \implies backups' = backupCount \\
  shared' = shared \\
  exclusive' = exclusive \\
  maxBytes' = maxBytes \\
  backupCount' = backupCount
\end{schema}

% Under LOCK_EX, the re-check found the file already small: a peer rotator ran
% first. The rotate is a no-op -- idempotency that closes the double-rotate race.
\begin{schema}{SkipRotate}
  \Delta Rotator \\
  w?: WRITER
\where
  w? \in exclusive \\
  size \leq maxBytes \\
  size' = size \\
  backups' = backups \\
  shared' = shared \\
  exclusive' = exclusive \\
  maxBytes' = maxBytes \\
  backupCount' = backupCount
\end{schema}

% The rotator releases LOCK_EX; appenders may re-acquire LOCK_SH and continue.
\begin{schema}{EndRotate}
  \Delta Rotator \\
  w?: WRITER
\where
  w? \in exclusive \\
  exclusive' = \emptyset \\
  shared' = shared \\
  size' = size \\
  backups' = backups \\
  maxBytes' = maxBytes \\
  backupCount' = backupCount
\end{schema}

\end{document}
